% Vérifié : citations, dates et renvois contrôlés sur lecture/p1-p2 et codicologie.yaml ; « velleités » rétabli tel qu'écrit ; visite Layral, Nazareth et bise du 31 précisées ; « perdu » (Journal sur feuille spéciale), « crise des tramways » (non sourcée) et rapprochement charbon/compartiments surchauffés retirés ; durée du retour, règle des trois mois, paie (78 f 75 puis 43 f 75) et billet à 11 f 15 précisés ; décès familiaux de l'introduction biographique allégés et donnés pour tradition ; « n'y reparaîtra plus » ramené à « n'y est plus nommée ».
% chapitre_VI.tex -- introduction du chapitre VI (clichés 128-146 ; ff. 64 r°-73 r°), fragment LaTeX sans préambule.
% Sources : carnet (lecture/p1), commentaires historiques (commentaire/128-146), dossiers R6, R1, R3, R4, R8.
% Rédaction du 18 oct. 2026.

Le 25~janvier 1917 au soir, Hippolyte Combes est à Acheux-en-Amiénois depuis bientôt trois mois. Logé chez un tailleur du village, les Carluvant (f.~51~r\textsuperscript{o}), il sert d'interprète au service d'achat de paille et de foin de la Fifth Army, sous les ordres du capitaine Williams, qui a remplacé le lieutenant Curtis le 2~décembre (f.~52~v\textsuperscript{o}). La 18th Division est revenue le 16~janvier sur l'Ancre, au sud de la rivière ; le V Corps a réinstallé le 12 son quartier général à Acheux, et c'est au bureau de la Mission française de ce corps que Combes a appris du capitaine Huguet la date de son congé, le 13 : «~mon départ fixé definit au 26 ct~» (f.~61~v\textsuperscript{o})\footnote{G.~H.~F.~Nichols, \textit{The 18th Division in the Great War}, Édimbourg et Londres, Blackwood, 1922, p.~136 et 139 ; C.~Falls, \textit{Military Operations, France and Belgium, 1917}, vol.~I, Londres, Macmillan, 1940 (ci-après \textit{OH 1917}, I), p.~66--69 ; C.~Baker, \textit{The Long, Long Trail}, «~The British V Corps in France and Flanders~» (en ligne).}. La veille du départ, par un gel intense, le carnet note encore : «~Nuit \& matinée, violent bombardem\textsuperscript{t}~» (f.~63~v\textsuperscript{o}). Les treize jours qui suivent sont la seconde et dernière parenthèse civile du \textit{Livre-Journal} ; à la différence de la première, elle suit une crise des fiançailles et conduit Combes, en fiancé, chez les parents de Marguerite Calmels, à Millau, sa ville natale. Au retour, le service qu'il avait laissé à Acheux n'y est plus.

\subsection*{Le carnet}

Le voyage d'aller prend près de deux jours. Le vendredi 26, Combes fait porter ses deux sacs à la gare d'Acheux-Varennes par la voiture de la poste, puis s'y rend à pied ; dans le compartiment de seconde classe, il décrit les champs couverts d'une neige durcie qui étincelle «~cme des grains de diam~» et échange avec un autre interprète «~nos impressions \& jugements sur les English~» (f.~64~r\textsuperscript{o} ; cf. \textit{Feuillets de Route}, chap.~\textsc{xvi}--\textsc{xviii}, annexe, p.~\pageref{annexe:doc6}). À Gézaincourt, il attend six heures dans une salle «~surchauff~» pleine de soldats permissionnaires et de paysans ; le train de Doullens, en retard de trois heures, le dépose à Amiens vers deux heures du matin. Avec deux collègues, dont Boilly, «~prof. à Pont-à-Mousson~», il trouve complet l'abri du soldat du boulevard d'Alsace-Lorraine, se réfugie au foyer de la Y.M.C.A. voisin de la gare et attend sur le quai le train de Boulogne, arrivé vers 5~h~30. Arrivé gare du Nord à 11~h~15, il fait enregistrer ses bagages à la gare d'Orsay (f.~65~r\textsuperscript{o}). Il déjeune à Neuilly, rue d'Armenonville, chez Mme Faure, avec «~Mlle Germaine \& Marcel~» ; il passe ensuite boulevard Saint-Michel chez l'opticien Imbert, «~mobilisé~», et fait quelques achats aux magasins de l'Hôtel-de-Ville avant de reprendre à Orsay, vers 19~h~50, le train de Toulouse (f.~65~v\textsuperscript{o}).

Le séjour toulousain va du dimanche 28~janvier au samedi 3~février. En approchant de la ville, les «~souv\textsuperscript{rs} jeunesse~» lui reviennent «~en foule~» ; il croit revoir sa mère et ses sœurs l'attendant à la gare, «~cme autrefois~». Il range ses affaires dans son appartement de la rue de Languedoc et monte à midi à Sainte-Marie, la maison des Crouzet sur la Côte-Pavée (f.~66~r\textsuperscript{o}), où Marguerite l'attend «~sur la porte~» ; il y passe la journée, en «~Causeries délicieuses sur les sujets les plus divers~», et redescend à pied, faute de tram (f.~66~v\textsuperscript{o}). Le 29, il retire en gare, «~non sans peine~», un panier de fleurs venu de Nice, puis visite avec Marguerite et Madeleine la basilique Saint-Sernin et ses cryptes : «~Toute ma jeunesse réveillée~» (f.~67~r\textsuperscript{o}). Le 30, après une visite à Mme de Labarthe, il assiste à la leçon publique d'Ernest Zyromski sur Maurice de Guérin, «~Vrai regal littéraire~» (f.~67~v\textsuperscript{o}). Le 31, sous une «~bise glaciale~», les fiancés font ensemble le marché de la halle Victor-Hugo, puis visitent l'hôtel d'Assézat, la Daurade et la Dalbade (f.~68~r\textsuperscript{o}). Le 1\textsuperscript{er}~février, il passe à Nazareth, où il a «~revu t\textsuperscript{te} ma famille~», choisit une montre en or chez le bijoutier Pfister et se rend à la chapelle de l'archevêché pour une conférence qui n'a «~pas eu lieu~». Le soir, la ville lui paraît silencieuse : «~Je me sens complet\textsuperscript{t} détaché de Toulouse. Partirai sans regret, et pour n'y retourner jamais plus~» (f.~68~v\textsuperscript{o}). Le lendemain pourtant, entre deux séances chez le dentiste et l'achat d'un «~remontoir or mat Louis XVI~» de 105~fr., il regrette de sacrifier «~3 jours entiers~» à Millau au lieu de jouir de son appartement, «~rempli de souvenirs~» (f.~69~r\textsuperscript{o}).

Le 3~février, Combes quitte Sainte-Marie avec Marguerite ; l'oncle les accompagne à la gare. À Carcassonne, il embrasse Jean, sans doute l'abbé Jean Crouzet, fils de l'oncle Émile ; à Béziers, entre deux trains, ils montent à la cathédrale, dont la terrasse découvre «~Plaine, rivière et soleil couchant~». Le «~petit train~» de la ligne des Causses, aux «~Vieux wagons~», les porte à Millau à 22~h~30 ; Pierre, vraisemblablement Pierre Calmels, frère de Marguerite, les attend, et Combes note l'«~Impression silence profond en debarquant ville natale~» (f.~69~v\textsuperscript{o}). La première nuit dans la maison Calmels est une nuit de deuil : le souvenir de «~Papa, Maman, Soeurs~» le fait sangloter, et il écoute sonner les heures jusqu'au matin (f.~70~v\textsuperscript{o}). Le dimanche, Marguerite lui joue quelques morceaux et chante ; il rend visite, «~cedant aux instances de tous~», à «~l'oncle \& tante Therèse (Layral)~», et prolonge la soirée avec Mimi, sans doute une autre sœur de Marguerite, arrivée par le train de 18~h. Pour les cloches du matin, il renvoie à un autre «~Journal~», tenu sur «~feuille speciale~» (f.~71~r\textsuperscript{o} ; cf. le feuillet de division des \textit{Feuillets de Route}, p.~\pageref{annexe:doc6}). Le 5, les fiancés vont seuls au cimetière ; les «~Adieux, dechirants~» se font vers cinq heures, et le train venu de Béziers, parti avec un long retard, s'arrête à toutes les stations sous la neige jusqu'à Saint-Flour, atteint à minuit et demi (f.~71~v\textsuperscript{o}).

Le retour prend plus de quarante heures, dont deux nuits de voyage. À Saint-Flour, Combes passe la nuit dans la salle d'attente, puis dans un compartiment de première classe garé sur la voie, les «~Pieds gelés~» ; à Brioude, il prend l'express P.L.M. de Paris, par Clermont, Gannat, Gien et Montargis, sous une neige continue (f.~72~r\textsuperscript{o}). À Paris, il gagne en métro la gare du Nord à temps pour le train de 19~h~19, plein de soldats «~fort curieux à observer~» (f.~72~v\textsuperscript{o}). D'Amiens, où il arrive à minuit et demi et se repose une heure au foyer de la Y.M.C.A., le train de Doullens le ramène lentement vers l'Ancre. En gare de Louvencourt, deux avions allemands évoluent au-dessus de la voie au milieu des éclatements : la guerre reparaît dans le carnet. Il est à Acheux vers 13~h~30, où les Carluvant lui offrent un «~déjeuner sommaire~» (f.~73~r\textsuperscript{o}).

\subsection*{Le contexte}

Les dates de ce congé s'accordent avec la règle commune. En mars 1916, la permission était «~uniformément de six jours pleins~», trajet non compris, et supposait trois mois de présence aux armées depuis la précédente\footnote{\textit{Bulletin des armées de la République}, 2\textsuperscript{e}~année, n\textsuperscript{o}~186, 22~mars 1916, p.~2 (La contemporaine, Argonnaute).} ; à partir du 1\textsuperscript{er}~octobre 1916, les combattants français eurent droit à trois permissions annuelles de sept jours\footnote{E.~Cronier, \textit{Permissionnaires dans la Grande Guerre}, Paris, Belin, 2013, p.~18 et 32, cité d'après l'article «~Permissionnaires français de la Première Guerre mondiale~» de Wikipédia (consulté en octobre 2026) : référence de seconde main.}. Combes repart quatre mois, à un jour près, après son départ de septembre 1916. L'interprète n'est pas payé moins pendant l'absence : il touche 78~f~75 de la Mission à son retour à Acheux (f.~74~r\textsuperscript{o}), et la «~retenue de 35 fr. pour 7 jours de permission~» n'est faite que le 3~mars, sur un versement réduit à 43~f~75 (f.~81~r\textsuperscript{o})\footnote{Sur la composition de la paie et la retenue de 5~F par jour de permission, voir l'introduction, «~Un interprète français dans l'armée britannique~».}. Le voyage, payé à tarif réduit, au «~quart de place~» dû à tout militaire en permission (le billet de Paris à Toulouse coûte 11~f~15, f.~66~r\textsuperscript{o}), se fait dans des trains surchargés : les cinq grands réseaux firent circuler 37\,674 trains de permissionnaires en 1916 et 81\,325 en 1917\footnote{G.~Lafon, \textit{Les Chemins de fer français pendant la guerre}, thèse de droit, Paris, A.~Rousseau, 1922, p.~37 (article~54 des cahiers des charges) et p.~152--153.}.

Le congé tombe au cœur de la vague de froid qui, du 16~janvier au 16~février environ, s'étend à toute la France : −18,5~°C à Clermont-Ferrand le 30~janvier, −15,4~°C à Paris et −8,3~°C à Toulouse le 4~février\footnote{G.~Barbé, «~L'hiver 1916-1917~», \textit{Journal d'agriculture pratique}, n.~s., t.~30, Paris, Librairie agricole de la Maison rustique, 1917, p.~181--182.}. La Seine charrie des glaçons depuis le 29~janvier ; à Paris, la crise du charbon est telle que le 27, jour du passage de Combes, le gouverneur militaire, le général Dubail, met voitures et hommes de corvée à la disposition des municipalités\footnote{\textit{La Dépêche} (Toulouse), 28~janv. 1917, p.~2, et 30~janv. 1917, p.~3 (Gallica et Rosalis).}. Les compartiments où Combes étouffe, à l'«~air brûlant, irrespirable~» (f.~65~r\textsuperscript{o}), contrastent avec cette pénurie. À Toulouse, la Faculté des lettres maintient ses cours publics : la leçon de Zyromski du 30~janvier, «~M. de Guérin : Le séjour en Bretagne et l'influence de Lamennais~», est annoncée puis résumée dans les journaux\footnote{\textit{Le Télégramme} (Toulouse), 30~janv. 1917, p.~3 ; \textit{L'Express du Midi}, 1\textsuperscript{er}~févr. 1917, p.~3, «~Maurice de Guérin à la Chênaie~» ; E.~Zyromski, \textit{Maurice de Guérin}, Paris, A.~Colin, 1921, chap.~VI, p.~113--114.}.

Sur l'Ancre, le gel, commencé vers le 14~janvier, a durci le sol et le thermomètre est descendu à −9,4~°C le 25. La 63rd Division enlève \textit{Puisieux Trench} et \textit{River Trench} du 3 au 5~février, l'évacuation de Grandcourt est constatée le 6, la ferme de Baillescourt occupée le 7, jour du retour de Combes\footnote{\textit{OH 1917}, I, p.~66--72 et 76.}. Le 4~février, Guillaume~II a signé l'ordre de l'opération \textit{Alberich}, le repli sur la ligne Hindenburg\footnote{\textit{Ibid.}, p.~110--115.} ; le 1\textsuperscript{er}, l'Allemagne a ouvert la guerre sous-marine à outrance, et les États-Unis ont rompu le 3 leurs relations diplomatiques avec elle\footnote{W.~Miles, \textit{Military Operations, France and Belgium, 1916}, vol.~II, Londres, Macmillan, 1938, p.~531, n.~4.}. Aucun de ces faits n'entre dans le carnet, qui ne retrouve la guerre qu'avec les deux «~taubes~» de Louvencourt.

Dans la vie de Combes, la permission efface la crise de décembre, où il avait noté des «~velleités de rupture~» (f.~56~v\textsuperscript{o} ; voir l'introduction biographique). Le séjour de Millau le fait entrer dans la parenté du docteur Gabriel Calmels, médecin, qui avait bâti la maison du 10, rue Alfred-Guibert : frère et sœurs de Marguerite, Layral, oncle et tante maternels, enfants Vincens ; Mme Faure, à Neuilly, serait une cousine issue de germain de la jeune fille\footnote{P.~Combes, \textit{Une étude de la famille Combes de Castelnau-Pégayrols}, dactylographie, version du 19~mars 2023, annexe~2, «~Les Calmels~», et étude familiale (arbres), 2020--2021, p.~106--107, archives familiales : tradition familiale, non recoupée par l'état civil sauf pour la profession du D\textsuperscript{r}~Calmels (acte de naissance de Marguerite Calmels, Millau, 16~avril 1892).}. Millau est aussi la ville de ses morts : le caveau familial, où le corps de Joséphine, morte à Nogent-le-Rotrou en octobre 1914, aurait été ramené selon la tradition (voir l'introduction biographique), est sans doute la tombe que les fiancés vont voir au cimetière le 5\footnote{P.~Combes, «~Mon grand-père Hippolyte Combes (et ses sœurs)~», dactylographie, s.~d., archives familiales.}. Le remontoir de 105~fr., avec le sautoir de 295~f acheté chez le même bijoutier en septembre (f.~28~v\textsuperscript{o}), représente environ deux mois et demi des versements de la Mission, qui vont de 125 à 165~F par mois ; l'encaisse notée au départ, 319~f~50 (f.~64~v\textsuperscript{o}), en valait à peu près deux\footnote{Comptes du carnet : sautoir, f.~28~v\textsuperscript{o} ; remontoir, f.~69~r\textsuperscript{o} ; sur les versements de la Mission, voir l'introduction, «~Un interprète français dans l'armée britannique~».}. Le mariage sera célébré à Millau le 25~octobre 1917, au cours d'une autre permission\footnote{Menu du repas, «~Menu --- 25~octobre 1917~», et photographie de mariage, archives familiales.}.

Le 7~février, en se rendant au bureau d'Acheux, Combes apprend que le capitaine Williams est parti pour Hérissart ; il fait ses paquets le jour même et quitte «~avec 1 serrement de cœur~» la maison des Carluvant (f.~73~v\textsuperscript{o}). Le chapitre suivant s'ouvre dans ce village, à la veille du dégel et du repli allemand ; Marguerite, si présente dans ces pages, n'y est plus nommée jusqu'à la fin du carnet.
